\documentclass[12pt,a4paper,twocolumn]{article}
\usepackage[margin=0.85in]{geometry}
\usepackage{times}
\usepackage{graphicx}
\usepackage{booktabs}
\usepackage{amsmath,amssymb}
\usepackage{xcolor}
\usepackage{hyperref}
\usepackage{enumitem}
\setlength{\parindent}{0.3in}

\begin{document}

% ------------------------- Front matter (spans both columns) ---------------
\twocolumn[
\begin{center}
{\LARGE\bfseries Climate-Smart Crop Calendar Recommendation and Monsoon\\
Shift Analysis for Tamil Nadu Using Machine Learning}
\end{center}

\vspace{6pt}

\begin{center}
{\bfseries Narmatha GI\textsuperscript{1}, Dr. Anita Sofia V S\textsuperscript{2}}\\
\vspace{4pt}
{\itshape \textsuperscript{1}PG Student, Department of Computer Application (PG), PSG College of Arts and Science, Coimbatore, Tamil Nadu, India.}\\
\url{25mca033@psgcas.ac.in}\textsuperscript{1}\\
\vspace{2pt}
{\itshape \textsuperscript{2}Associate Professor, Department of Computer Application (PG), PSG College of Arts and Science, Coimbatore, Tamil Nadu, India.}\\
\url{anithasofia@psgcas.ac.in}\textsuperscript{2}
\end{center}

\vspace{10pt}

\noindent\textbf{ABSTRACT}

\noindent Rainfall in Tamil Nadu is not static from year to year, yet most machine-learning
crop recommendation systems treat it as a single current-season feature rather than
asking whether a district's monsoon behavior has itself been shifting over the
available historical record. This study presents a climate-smart crop calendar
system for Tamil Nadu that is built around two coupled components: a district-level
monsoon shift analysis over each district's real 21-year historical rainfall record
(2005--2025), and a machine-learning yield and sowing-window recommendation pipeline
conditioned on that district's own climate history rather than a generic national
average. The system draws on real government agricultural census data (38 districts,
25 crops, six government-defined cropping seasons, 2005--2019) merged with
district-level historical climate data, combined with a TNAU-style agronomic
sowing-window reference calendar. While preparing this system, a systematic
$\sim$100$\times$ production-unit inconsistency was discovered in the underlying
government yield records for every crop in the dataset, silently corrupting a
substantial share of the training signal; an empirical, duplicate-record-grounded
correction procedure is described (6,295 affected records rescaled, one unrecoverable
record dropped). On the corrected data, three regression models were trained and
fairly compared for district-level yield estimation -- Decision Tree, Random Forest,
and XGBoost -- with hyperparameters tuned under matched trees$\times$depth budgets.
The tuned XGBoost model achieved the best held-out performance (test $R^2 = 96.98\%$,
MAE $= 1.45$, RMSE $= 3.82$) against Random Forest ($R^2 = 96.57\%$) and a single
Decision Tree ($R^2 = 95.27\%$), confirmed with 5-fold cross-validation. The monsoon
shift analysis found 4 of 38 districts with a statistically non-stable rainfall trend
over the analysed period, including declines of 46.55\% (Tenkasi) and 41.90\%
(Tirunelveli), each of which shifts that district's drought-risk classification
independently of the current season's rainfall. The resulting five-stage
recommendation pipeline is deployed as a full-stack, multilingual web application
(Flask, PostgreSQL) with live weather integration, installable as both a Progressive
Web App and a native Android application.

\vspace{6pt}
\noindent\textbf{Keywords:} Crop Recommendation, Monsoon Shift Analysis, Rainfall
Trend, Machine Learning, XGBoost, Precision Agriculture, Data Quality, Deployed
Decision Support System

\vspace{10pt}
]

\section*{I. INTRODUCTION}

Machine learning (ML) has become a standard tool for crop recommendation and yield
prediction, supporting farmers with data-driven guidance on what to grow and when
[1]. A large share of published crop recommendation studies, however, train and
validate models on a single widely-used public dataset of 2,200 samples spanning 22
crop classes with seven generic agro-climatic features (nitrogen, phosphorus,
potassium, temperature, humidity, pH, and rainfall) [1]. This dataset is convenient
and well-balanced, but it is not tied to any specific region's actual administrative
and agronomic structure -- it has no district, no government-defined cropping season,
and no historical rainfall trend for a real place a farmer actually farms in.
Consequently, models trained on it, however accurate on their own test split, cannot
directly tell a farmer in a named district when to sow a named crop this season, and
cannot say whether that district's monsoon has itself been changing.

A second, related gap is that even studies which do use real rainfall data typically
treat it as one more static input feature at prediction time, rather than first
asking whether the region's monsoon behavior itself has been \emph{shifting} over the
available historical record. Trend-detection studies of Indian rainfall [16], by
contrast, show that this question has a real, non-trivial answer -- most Indian
meteorological divisions exhibit statistically significant long-term rainfall trends
-- and that answer is exactly the kind of district-specific signal a crop calendar
recommendation should be conditioned on, not just a raw rainfall number for the
current season.

This study takes both gaps as its starting point. Instead of a generic feature-based
dataset and a static rainfall feature, a crop recommendation and sowing-calendar
system is built directly from government agricultural census and district-level
climate records for the Indian state of Tamil Nadu, each district's real multi-decade
rainfall trend is explicitly analyzed as a monsoon-shift signal that feeds the
recommendation's risk assessment, and the resulting model is carried through to a
deployed, installable application rather than remaining an offline evaluation. This
process surfaced a data problem that a small curated benchmark dataset would not:
real government production records, aggregated over decades from multiple reporting
cycles, contain unresolved unit inconsistencies that directly affect the validity of
any yield model -- and any rainfall-trend analysis -- trained on this class of data
source. This issue and its correction procedure are detailed as one of this study's
contributions.

The contributions of this study are:
\begin{enumerate}[label=(\roman*)]
  \item A district-level monsoon shift analysis over each Tamil Nadu district's real
        21-year historical rainfall record, translating the detected trend direction
        and magnitude into an explicit drought/flood risk signal that conditions the
        downstream crop recommendation, rather than treating rainfall as a single
        static feature.
  \item A district-level crop sowing recommendation pipeline built on real Tamil Nadu
        government agriculture and climate records (25 crops, 38 districts, six
        government-defined seasons), combining a TNAU-style agronomic reference
        calendar with the monsoon-shift-based risk assessment above.
  \item Identification and correction of a systematic $\sim$100$\times$
        production-unit inconsistency affecting every crop in the underlying
        government dataset, using an empirical procedure grounded in confirmed
        duplicate-record pairs rather than a purely statistical outlier heuristic.
  \item A fair, matched-budget hyperparameter comparison of Decision Tree, Random
        Forest, and XGBoost regressors for yield estimation on the corrected data,
        validated with both a held-out split and 5-fold cross-validation.
  \item A deployed, farmer-facing full-stack implementation (Flask, PostgreSQL)
        installable as a Progressive Web App and as a native Android application,
        with live weather data replacing static estimates on the dashboard.
\end{enumerate}

\section*{II. RELATED WORK}

Extensive research has investigated machine learning approaches for crop
recommendation, yield prediction, and rainfall trend analysis. This section reviews
contributions relevant to each of these three threads.

Alam et al. [1] proposed a soft-voting ensemble of Logistic Regression, Na\"ive
Bayes, KNN, Decision Tree, Random Forest, and XGBoost on a 2,200-sample generic crop
dataset, reporting 99.54\% classification accuracy and introducing entropy-based
uncertainty quantification over the ensemble's softmax probabilities. Hasan et al.
[2] used a K-nearest-neighbor/Random Forest/Ridge-regression ensemble for crop
production prediction in Bangladesh. Bandi et al. [3] and Nti et al. [4] each
reported classification accuracies above 99\% using voting and tree-based ensembles
respectively for crop recommendation. Gopi and Karthikeyan [5] combined Red Fox
Optimization with an ensemble recurrent neural network for joint crop recommendation
and yield prediction. Elbasi et al. [6] compared fifteen ML algorithms combined with
IoT sensor data, reporting 97\% accuracy. Bera et al. [7] proposed an edge/dew-
computing and federated-learning architecture (E-CropReco) to reduce latency in
real-time crop recommendation. Thilakarathne et al. [8] built a cloud-based IoT
platform for real-time crop monitoring, though scoped to a single crop. Shingade et
al. [9] developed a Random Forest classifier specifically for seed recommendation.
Senapaty et al. [10] and Chaudhari et al. [11] each describe decision-support-
oriented crop recommendation systems closer in spirit to a deployed tool than an
offline benchmark. Agrawal et al. [12] combined ML with spatial multi-criteria
decision analysis for agricultural land suitability classification, reasoning at the
level of a specific geographic unit rather than a generic feature vector. Afzal et
al. [13] combined Random Forest and XGBoost with soil information to improve
recommendation output.

Beyond crop recommendation, a closely related line of work addresses yield
prediction directly with tree-ensemble methods. Jeong et al. [14] evaluated Random
Forests against multiple linear regression for wheat, maize, and potato yield
prediction at global and regional scale, reporting substantially lower error for
Random Forest than the linear baseline. Hu et al. [15] cautioned against black-box ML
models for evaluating climate-change impacts on crop yield and argued for
explainable, interpretable approaches instead.

A separate line of work analyzes rainfall trends directly rather than using rainfall
only as a prediction feature. Praveen et al. [16] applied the Mann-Kendall trend
test, Pettitt change-point detection, and an ANN-MLP forecaster to 115 years
(1901--2015) of Indian meteorological-division rainfall data, finding statistically
significant negative rainfall trends in most divisions. Correndo et al. [17] combined
Bayesian statistics and machine learning to show that weather is the dominant source
of uncertainty in maize yield response to nitrogen. Rihan et al. [18] used machine
learning and deep learning with explicit sensitivity and uncertainty analysis for
forest-fire susceptibility mapping and identified rainfall as one of the dominant
driving factors. The foundational tree-ensemble algorithms used throughout this study
are Random Forest [19] and XGBoost [20].

These studies share two characteristics relevant to the present work. First, nearly
all are evaluated on generic or single-crop datasets rather than a real
administrative region's multi-decade agricultural census, so a trained model has no
direct notion of which named district, which government-defined season, and which
historical climate trend apply to a given farmer. Second, none combine crop
recommendation with an explicit, district-level rainfall-trend-based risk signal, and
none are carried through to a farmer-facing deployed web and mobile application. The
present study addresses both gaps directly.

\section*{III. METHODOLOGY}

\subsection*{A. System Workflow}

The proposed system implements an end-to-end pipeline that, given a district, crop,
and season, executes five stages: climate data retrieval, monsoon-shift-aware climate
risk assessment, yield prediction, sowing-window selection, and confidence scoring.
The workflow commences with data acquisition from real government agriculture and
climate records, proceeds through the data-quality correction described in Section
III-C, trains and compares three regression models (Section III-D), computes each
district's rainfall trend (Section III-E), and serves the resulting recommendation
through a deployed web and mobile application (Section III-H).

\subsection*{B. Data Collection and Preprocessing}

The system is built on two merged real data sources: a government agricultural
production census covering Tamil Nadu, reporting district, year, season, crop, area
under cultivation, and production for 2005--2019; and district-level historical
climate records (temperature, rainfall, humidity, wind) spanning 2005--2025 for the
same districts. The two sources are merged on district, year, and season to produce a
single feature table used both for model training and the application's real-time
climate-retrieval logic. The application supports all 38 present-day Tamil Nadu
districts; five districts created after the 2019 administrative reorganization
(Kallakurichi, Mayiladuthurai, Ranipet, Tenkasi, Tirupathur) are mapped to the parent
district they were carved from for historical lookups, while remaining independently
selectable in the interface.

\subsection*{C. Data Preprocessing Techniques}

To address a data-sufficiency gap and a data-quality defect discovered in the raw
government source, three preprocessing techniques were implemented.

\textbf{Crop Selection Criterion:} The government census records production for many
more crops than the eight originally covered by an early version of this system. To
responsibly expand coverage, every additional crop reported for Tamil Nadu was
included only if it met all three of: at least 300 recorded instances, at least 15
distinct years of history, and at least 15 distinct districts. Seventeen crops
cleared this bar and were added, bringing total coverage to 25 crops. Crops that
failed the bar (e.g., Mango, Grapes, and Papaya, each with only two recorded years)
were deliberately excluded rather than included with unreliable history; Tea, Coffee,
and Rubber are not present in this census source at all, since they are tracked by
separate commodity boards in India.

\textbf{Production-Unit Correction:} Yield figures for individual (district, crop,
year) combinations were found to be physically implausible for a subset of records --
e.g., rice in one district reporting approximately 3--4 tonnes/hectare for
2005--2013 and approximately 330--490 tonnes/hectare for 2014--2019, an impossible
jump for the same crop and district. Manual inspection revealed that many records are
recorded twice, with one copy's Production value exactly 100$\times$ the other,
confirmed on thousands of matching duplicate-row pairs. An existing deduplication
step already preferred the smaller value when both copies existed, but a large share
of records had only the inflated copy submitted, leaving nothing to prefer. This
bimodal pattern was confirmed present in every one of the 25 crops. A purely
statistical correction (largest gap in the sorted value distribution) was evaluated
and rejected because it let some crops' genuinely implausible values through as
legitimate. Instead, a correction was derived empirically: for each crop, the 99th
percentile of the correct-scale member of every confirmed duplicate pair (ratio
90--110$\times$) was computed, and each crop's correction ceiling was set to
1.5$\times$ this percentile. Records exceeding their crop's ceiling had Production
divided by 100; records still more than 3$\times$ their ceiling afterward were
dropped rather than further guessed at. In total, 6,295 records were rescaled and one
record was dropped (a Cashewnut observation with an unrelated broken Area value of
1.0 hectare) across the full 25-crop dataset.

\textbf{Per-Crop Yield Categorization:} The Low/Medium/High yield-tertile label used
elsewhere in the pipeline was found to exhibit a ``crop-identity leak'': because raw
yield scales differ enormously by crop, a single global quantile split classified
nearly every high-yield crop as ``High'' and every low-yield crop as ``Low''
regardless of how each performed relative to its own crop's normal range. This was
corrected by computing yield quantiles per crop rather than globally.

\subsection*{D. Model Architecture -- Regression Models}

Three regression models were trained on the corrected 25-crop feature table with an
80:20 train/test split: a single Decision Tree regressor as a baseline; a Random
Forest regressor [19] (\texttt{n\_estimators=300}, \texttt{max\_depth=20},
\texttt{min\_samples\_split=5}, \texttt{min\_samples\_leaf=2}); and an XGBoost
regressor [20]. Twenty features are used: district, year, season, crop, cultivated
area, six raw climate statistics, three derived statistics (rainy days, temperature
and rainfall standard deviation), and four categorical bucket features (temperature,
rainfall, humidity, wind).

An initial XGBoost configuration used substantially more conservative hyperparameters
than the Random Forest baseline (\texttt{max\_depth=3}, strong L1/L2
regularization), which is not a fair comparison, and under-performed Random Forest as
a result. XGBoost was therefore re-tuned to a comparable trees$\times$depth budget
(\texttt{n\_estimators=800}, \texttt{max\_depth=6}, \texttt{learning\_rate=0.03},
\texttt{min\_child\_weight=1}, \texttt{subsample=0.85},
\texttt{colsample\_bytree=0.85}, \texttt{reg\_alpha=0.01}, \texttt{reg\_lambda=1.0}),
selected by checking both held-out test $R^2$ and 5-fold cross-validated $R^2$
together.

\subsection*{E. Monsoon Shift Analysis}

For each of the 38 districts, a 21-year rainfall trend (2005--2025) is computed from
the historical climate record: percentage change over the period and a trend
direction/strength classification. This mirrors the trend-then-risk pattern used in
Indian rainfall trend-detection studies such as [16], though computed per-district
here rather than at the coarser meteorological-division scale. Drought and flood risk
for a recommendation are derived jointly from the district's current-season rainfall
category \emph{and} this historical trend, rather than from current-season rainfall
alone -- so a district with currently-adequate rainfall but a strong long-term
declining trend is flagged at elevated drought risk rather than ``Low.''

\subsection*{F. Supporting Pipeline Stages}

\textbf{Sowing Window Selection:} A TNAU-style (Tamil Nadu Agricultural University)
reference calendar maps each (crop, season) pair to an agronomic sowing-window date
range. Because the underlying dataset contains no sowing-date ground truth, this
stage relies on domain-expert reference values. If the current year's window has
already elapsed, the system advances the recommendation to the next occurrence and
surfaces this explicitly to the farmer rather than silently recommending a passed
date.

\textbf{Confidence Scoring:} A data-driven confidence score is derived from the
predicted yield's per-crop tertile category and the overall drought/flood risk level.

\subsection*{G. Model Training Strategy}

Models were trained with an 80:20 train/test split (\texttt{random\_state=42}).
Model selection was validated with 5-fold cross-validation in addition to the
held-out split, to confirm that observed differences between models reflected real
generalization behavior rather than the particular random split chosen.

\subsection*{H. System Architecture}

The trained model is served through a full-stack web application. The backend
(Python, Flask) exposes the five-stage pipeline as a farmer-facing dashboard, a
district/crop/season recommendation form, a monsoon/climate-trend view, a sowing
calendar view, and a notification center, in English and Tamil. Application data is
persisted in PostgreSQL. The dashboard's live weather figures are retrieved per
district from the Open-Meteo API at request time, replacing an earlier static
per-district estimate. The same application is exposed as an installable Progressive
Web App (web manifest and service worker) and as a native Android application built
with Capacitor, wrapping the same served pages in a WebView shell rather than
re-implementing the interface natively. The system is deployed to a cloud host with a
managed PostgreSQL instance, reachable from any network.

\section*{IV. DATASET DESCRIPTION}

The system uses real Tamil Nadu government agricultural census records (2005--2019)
merged with district-level historical climate records (2005--2025). Coverage spans
38 districts, 25 crops, and six government-defined cropping seasons (Kharif, Rabi,
Summer, Autumn, Winter, Whole Year). The 25 crops comprise the original eight (Rice,
Groundnut, Cotton, Sugarcane, Maize, Bajra, Ragi, Banana) plus seventeen added under
the criterion in Section III-C: five pulses (Urad, Moong, Horse-gram, Arhar/Tur,
Gram), two oilseeds (Sesamum, Sunflower), two millets (Jowar, Small millets), three
spices (Turmeric, Dry chillies, Coriander), one plantation crop (Cashewnut), and four
others (Tobacco, Onion, Tapioca, Sweet potato). The production-unit correction
(Section III-C) rescaled 6,295 records and dropped one unrecoverable record across
this dataset.

The monsoon shift analysis (Section III-E) computed a 21-year rainfall trend for all
38 districts. Of these, 34 districts were classified as having a stable long-term
rainfall pattern, while 4 districts showed a statistically non-stable trend --
including a 46.55\% decline in Tenkasi and a 41.90\% decline in Tirunelveli over the
analysed period, the two most pronounced shifts detected.

\section*{V. EXPERIMENTAL SETUP}

\subsection*{A. Hardware and Software Configuration}

Unlike deep-learning computer-vision pipelines, gradient-boosted and bagged tree
models of this scale do not require GPU acceleration; all training was performed on
CPU. The software stack comprises Python for the modeling pipeline; pandas, NumPy,
and scikit-learn for data preparation and the Decision Tree/Random Forest models;
XGBoost for the gradient-boosted model; joblib for model persistence; Matplotlib for
result visualization; Flask for the backend web application; PostgreSQL (via
psycopg2) for persistence; and, for the mobile targets, a Node.js/Capacitor toolchain
with the Android SDK and Gradle to build the native Android package. The application
is deployed to a managed cloud host running Python via Gunicorn, with a managed
PostgreSQL instance.

\subsection*{B. Model Training and Evaluation}

Training used an 80:20 train/test split with \texttt{random\_state=42} for
reproducibility. Evaluation metrics were $R^2$ (train and test), Mean Absolute Error
(MAE), and Root Mean Squared Error (RMSE), together with the train-test $R^2$ gap as
an overfitting indicator. Five-fold cross-validation was additionally run to confirm
that the comparison between models was not an artifact of the particular 80:20 split.

\section*{VI. RESULTS AND DISCUSSIONS}

\subsection*{A. Regression Performance}

\begin{table*}[t]
\centering
\caption{Model Comparison on the Corrected Dataset (Held-Out Test Split)}
\begin{tabular}{lccccc}
\toprule
Model & Train $R^2$ & Test $R^2$ & MAE & RMSE & Gap \\
\midrule
XGBoost & 99.60\% & \textbf{96.98\%} & \textbf{1.45} & \textbf{3.82} & 0.026 \\
Random Forest & 98.38\% & 96.57\% & 1.19 & 4.08 & 0.018 \\
Decision Tree & 97.60\% & 95.27\% & 1.46 & 4.79 & 0.023 \\
\bottomrule
\end{tabular}
\end{table*}

\noindent XGBoost achieved the best test $R^2$ (96.98\%), ahead of Random Forest
(96.57\%) and the single Decision Tree (95.27\%), with the lowest MAE and RMSE among
the three. Five-fold cross-validation confirmed the result: the tuned XGBoost
configuration achieved a mean CV $R^2$ of 0.9585 versus Random Forest's 0.9576 and
the original, under-tuned XGBoost configuration's 0.9467, showing that fair tuning,
not the particular split, drives the improvement.

\subsection*{B. Overfitting Gap Analysis}

Random Forest has the lowest train-test $R^2$ gap of the three models (0.018 versus
XGBoost's 0.026). Both gaps, however, are comfortably inside a ``good
generalization'' threshold (gap $\leq 0.05$); neither model shows concerning
overfitting. Overfitting gap measures the consistency of a model's score between
training and test data, not its absolute quality -- a model that always predicted the
per-crop mean yield would have a gap near zero while being useless. Once both models
are confirmed to generalize well, the appropriate tie-breaker is test-set accuracy,
which favors XGBoost. XGBoost's slightly larger gap is also structurally expected:
boosting fits training data more closely by design, whereas bagging's averaging
across independently-trained trees smooths training fit further -- a difference in
mechanism, not a flaw in either model.

\subsection*{C. Feature Importance Analysis}

Feature importance analysis of the trained XGBoost model shows climate variables
(rainfall and temperature statistics) and cultivated area dominating the prediction,
consistent with agronomic expectation, while the categorical bucket features
contribute comparatively little additional signal once their underlying continuous
values are already present.

\subsection*{D. Monsoon Shift Analysis Results}

Of the 38 districts analyzed, 34 (89\%) showed a statistically stable long-term
rainfall pattern over 2005--2025, while 4 districts showed a non-stable trend. The
two most pronounced declines were Tenkasi (--46.55\%) and Tirunelveli (--41.90\%),
both in southern Tamil Nadu. For these districts, the recommendation pipeline's
drought-risk classification is elevated independently of the current season's
rainfall reading, directly reflecting the contribution described in Section I.

\subsection*{E. Real-World Deployment}

The trained model is served through a full-stack web application rather than
remaining an offline artifact, installable as a Progressive Web App and as a native
Android application, with live per-district weather replacing an earlier static
estimate. This distinguishes the present system from most of the related work
reviewed in Section II, which stops at offline model evaluation.

\subsection*{F. Practical Implications}

The production-unit inconsistency described in Section III-C is, to the authors'
knowledge, not previously documented for this class of Indian government
agricultural census data, and has a direct implication for any prior or future model
trained on similar uncorrected sources: a substantial fraction of the yield target
variable may be corrupted by a consistent, detectable, but easily-overlooked unit
error. Grounding the correction in confirmed duplicate-record evidence, rather than a
purely statistical heuristic, is considered as significant a contribution of this
work as the model comparison itself.

\section*{VII. CONCLUSION AND FUTURE WORK}

\subsection*{A. Conclusion}

This study presented a district-level crop sowing recommendation system for Tamil
Nadu built around real government agricultural and climate records, coupled with an
explicit monsoon shift analysis over each district's 21-year rainfall history rather
than a single static rainfall feature. In the course of building it, a systematic
$\sim$100$\times$ production-unit inconsistency was identified and corrected in the
underlying government source for every crop, and a related per-crop yield-
categorization flaw was corrected. On the corrected data, a fairly-tuned XGBoost
regressor achieved the best held-out performance (test $R^2 = 96.98\%$) among
Decision Tree, Random Forest, and XGBoost baselines, verified with cross-validation.
The monsoon shift analysis identified four districts with statistically non-stable
rainfall trends, including declines exceeding 40\% in two southern districts. The
resulting pipeline is deployed as a full-stack, multilingual web application with
live weather integration, installable as both a Progressive Web App and a native
Android application.

\subsection*{B. Future Work}

Several directions remain for future work. Uncertainty quantification -- adapting an
entropy-based or ensemble-variance-based measure, in the style of [1], to this
regression setting -- would extend the coarse tertile-category and heuristic
confidence score currently used, particularly for districts or crops with sparse
historical records. The historical data underlying both the yield model and the
climate risk assessment currently ends in 2019 for agriculture and 2025 for climate;
incorporating more recent government releases as they become available, and
integrating real-time soil-sensor data, would improve currency. The sowing-window
calendar is currently a domain-expert reference rather than a learned quantity, since
the dataset contains no sowing-date ground truth; sourcing such records from
agricultural extension offices is a direct avenue to let the model learn
district-specific timing shifts. Finally, extending the same data-correction and
modeling methodology to other Indian states' government agricultural census data is a
direct avenue for broader impact.

\section*{VIII. REFERENCES}

\begin{enumerate}[label={[\arabic*]}, leftmargin=*, itemsep=4pt]
\item Alam, Md. S. B., Esichaikul, V., Lameesa, A., Ahmed, S. F., \& Gandomi, A. H.
(2025). An approach for crop recommendation with uncertainty quantification based on
machine learning for sustainable agricultural decision-making. \emph{Results in
Engineering}, 26, 105505.

\item Hasan, M., Marjan, M. A., Uddin, M. P., Afjal, M. I., Kardy, S., Ma, S., \&
Nam, Y. (2023). Ensemble machine learning-based recommendation system for effective
prediction of suitable agricultural crop cultivation. \emph{Frontiers in Plant
Science}, 14.

\item Bandi, R., Likhit, M. S. S., Reddy, S. R., Bodla, S. R., \& Venkat, V. S.
(2023). Voting classifier-based crop recommendation. \emph{SN Computer Science}, 4.

\item Nti, I. K., Zaman, A., Nyarko-Boateng, O., Adekoya, A. F., \& Keyeremeh, F.
(2023). A predictive analytics model for crop suitability and productivity with
tree-based ensemble learning. \emph{Decision Analytics Journal}, 8.

\item Gopi, P. S. S., \& Karthikeyan, M. (2024). Red fox optimization with ensemble
recurrent neural network for crop recommendation and yield prediction model.
\emph{Multimedia Tools and Applications}, 83.

\item Elbasi, E., Zaki, C., Topcu, A. E., Abdelbaki, W., Zreikat, A. I., Cina, E.,
Shdefat, A., \& Saker, L. (2023). Crop prediction model using machine learning
algorithms. \emph{Applied Sciences}, 13(16), 9288.

\item Bera, S., Dey, T., Mukherjee, A., \& Buyya, R. (2023). E-CropReco: a
dew-edge-based multi-parametric crop recommendation framework for internet of
agricultural things. \emph{The Journal of Supercomputing}, 79.

\item Thilakarathne, N. N., Bakar, M. S. A., Abas, P. E., \& Yassin, H. (2023).
Towards making the fields talk: a real-time cloud enabled IoT crop management
platform for smart agriculture. \emph{Frontiers in Plant Science}, 13.

\item Shingade, S. D., Mudhalwadkar, R. P., \& Masal, K. M. (2022). Random forest
machine learning classifier for seed recommendation. In \emph{Proceedings of the
International Conference on Edge Computing and Applications (ICECAA)}.

\item Senapaty, M. K., Ray, A., \& Padhy, N. (2024). A decision support system for
crop recommendation using machine learning classification algorithms.
\emph{Agriculture}, 14(8), 1256.

\item Chaudhari, R. R., Jain, S., \& Gupta, S. (2025). Agricultural machine learning
platform: enhancing crop suggestion and crop yield estimates. \emph{Journal of
Integrated Science and Technology}, 13.

\item Agrawal, N., Govil, H., \& Kumar, T. (2024). Agricultural land suitability
classification and crop suggestion using machine learning and spatial multicriteria
decision analysis in semi-arid ecosystem. \emph{Environment, Development and
Sustainability}.

\item Afzal, H., Amjad, M., Raza, A., Munir, K., Villar, S. G., Lopez, L. A. D., \&
Ashraf, I. (2025). Incorporating soil information with machine learning for crop
recommendation to improve agricultural output. \emph{Scientific Reports}, 15.

\item Jeong, J. H., Resop, J. P., Mueller, N. D., Fleisher, D. H., Yun, K., Butler,
E. E., Timlin, D. J., Shim, K.-M., Gerber, J. S., Reddy, V. R., \& Kim, S.-H. (2016).
Random forests for global and regional crop yield predictions. \emph{PLOS ONE},
11(6), e0156571.

\item Hu, T., Zhang, X., Bohrer, G., Liu, Y., Zhou, Y., Martin, J., Li, Y., \& Zhao,
K. (2023). Crop yield prediction via explainable AI and interpretable machine
learning: dangers of black box models for evaluating climate change impacts on crop
yield. \emph{Agricultural and Forest Meteorology}, 336, 109458.

\item Praveen, B., Talukdar, S., Shahfahad, Mahato, S., Mondal, J., Sharma, P.,
Islam, A. R. M. T., \& Rahman, A. (2020). Analyzing trend and forecasting of
rainfall changes in India using non-parametrical and machine learning approaches.
\emph{Scientific Reports}, 10, 10342.

\item Correndo, A. A., Tremblay, N., Coulter, J. A., Ruiz-Diaz, D., Franzen, D.,
Nafziger, E., Prasad, V., Rosso, L. H. M., Steinke, K., Du, J., Messina, C. D., \&
Ciampitti, I. A. (2021). Unraveling uncertainty drivers of the maize yield response
to nitrogen: a Bayesian and machine learning approach. \emph{Agricultural and Forest
Meteorology}, 311, 108668.

\item Rihan, M., Ali Bindajam, A., Talukdar, S., Shahfahad, Naikoo, M. W., Mallick,
J., \& Rahman, A. (2023). Forest fire susceptibility mapping with sensitivity and
uncertainty analysis using machine learning and deep learning algorithms.
\emph{Advances in Space Research}, 72.

\item Biau, G., \& Scornet, E. (2016). A random forest guided tour. \emph{Test}, 25.

\item Chen, T., \& Guestrin, C. (2016). XGBoost: a scalable tree boosting system. In
\emph{Proceedings of the ACM SIGKDD International Conference on Knowledge Discovery
and Data Mining}.

\end{enumerate}

\end{document}
